\section{Scope, Method, and Qualifications}

\subsection*{Audience, Purpose, and Series Placement}

This essay is written for intelligent adult readers---whether practicing Catholics seeking deeper doctrinal clarity or thoughtful inquirers engaging Catholic teaching from within its classical sources---without presupposing graduate training in scholastic terminology. It forms Part III of the trilogy \emph{The Creature Before God}. Whereas Part I examines received being and the Creator's sustaining presence, and Part II examines what justice owes to God and neighbor in worship, truth, and obedience, Part III investigates the ultimate end of human desire, the elevation of obedient justice into infused charity as friendship with God, the lived tension of hope, prayer, dryness, Eucharistic desire, suffering, and neighbor-love, and the consummation of creaturehood in the beatific vision.

Offered as a study aid in Catholic faith and theology, this work carries no magisterial authority, does not replace the official Catechism, liturgical books, or sacramental formation, and does not constitute medical, psychological, or individualized pastoral direction.

\subsection*{Epistemic Seam and Authority Boundaries}

Throughout the exposition, philosophical analysis accessible to natural reason and revealed truths received in divine faith remain explicitly distinguished. The inadequacy of finite goods to quiet a rational will, and the intellect's wonder at the First Cause, belong to philosophical inquiry; the actual ordination of humanity to the supernatural vision of the divine essence, the infusion of charity as friendship with God, the sacramental economy of the Eucharist, and the promise of the bodily resurrection belong strictly to divine revelation. Natural desire is never treated as a juridical claim or exigency binding God.

Doctrinal claims are stated at their received level of authority:
\begin{itemize}
\item \textbf{Sacred Scripture and solemn dogma:} the beatific vision of the divine essence (\emph{Benedictus Deus}, 1336), the real distinction between God and creation and the absolute necessity of revelation for our supernatural end (\emph{Dei Filius}, 1870), justification and the gratuity of first grace (Council of Trent, Session VI), and the real presence and conditions for worthy reception of the Eucharist (Council of Trent, Session XIII).
\item \textbf{Authoritative non-definitive teaching:} the \emph{Catechism of the Catholic Church} (1992/1997) and the encyclicals \emph{Ecclesia de Eucharistia} (St.\ John Paul II, 2003) and \emph{Spe salvi} (Pope Benedict XVI, 2007).
\item \textbf{Patristic, scholastic, and saintly doctors:} St.\ Irenaeus of Lyons, St.\ Augustine of Hippo, St.\ Bernard of Clairvaux, St.\ Thomas Aquinas, St.\ Teresa of Ávila, St.\ Francis de Sales, and St.\ Th\'er\`ese of Lisieux. Their arguments and images represent major received theology and spiritual teaching, without claiming that every scholastic formulation (such as the precise technical vocabulary of \emph{lumen gloriae} or Bernard's fourfold schema) carries the status of an ecumenical definition.
\end{itemize}

\subsection*{Pastoral and Moral Safeguards}

Treatment of spiritual dryness and human suffering is strictly bounded against misapplication. Spiritual dryness is distinguished from clinical depression or bodily illness, though the two can coincide; readers experiencing persistent psychological distress or suicidal thoughts require qualified medical and pastoral care. Suffering is affirmed as an evil, never a good in itself and never presumptive personal punishment; removable suffering ought to be resisted and relieved with lawful remedies; self-harm, passivity before injustice, and generalized severe physical austerity are expressly excluded. In sacramental matters, spiritual communion is presented alongside---never as a substitute for---sacramental Reconciliation in the case of conscious grave sin before receiving the Eucharist.

\Needspace{8\baselineskip}
\subsection*{Editions, Translations, Rights, and Negative Results}

All direct verbatim English Scripture quotations are verified against the tracked public-domain Challoner revision of the Douay-Rheims Bible (Project Gutenberg eBook 1581); scriptural phrases embedded inside block or inline quotations of Aquinas, Augustine, Bernard, or other authors retain the wording of the quoted translation. Quotations from St.\ Thomas Aquinas's \emph{Summa Theologiae} (Prima, Prima Secundae, Secunda Secundae, and Tertia Pars) are verified directly against the four tracked public-domain English Dominican translation texts (Benziger Brothers; Project Gutenberg eBooks 17611, 17897, 18755, and 19950). Patristic, conciliar, and spiritual-theological verbatim quotations are verified directly against tracked public-domain witnesses in the repository source library: St.\ Augustine's \emph{Confessions} (Pilkington, NPNF1-1 CCEL text) and \emph{City of God} (Dods 1871, Gutenberg eBook 45305, collated with Hoffmann's CSEL 40.2 Latin text); St.\ Irenaeus's \emph{Adversus Haereses} (Roberts--Rambaut, ANF-1 CCEL text); the First Vatican Council's \emph{Dei Filius} (Latin text collated with Cardinal Manning's 1912 Devin-Adair English text); St.\ Bernard of Clairvaux's \emph{On Loving God} (Marianne Caroline and Coventry Patmore, 1884); St.\ Teresa of Ávila's \emph{Life} (Lewis--Zimmerman, 1904); St.\ Francis de Sales's \emph{Introduction to a Devout Life} (undated Pustet New York edition, cross-checked against the 1832 P\'erisse Fr\`eres French text, Gutenberg eBook 53540); and St.\ Th\'er\`ese of Lisieux's \emph{Story of a Soul} (Burns \& Oates fifteenth-thousand edition edited by T.\ N.\ Taylor, 1912, strictly within the tracked autobiography segment). Minor optical-character-recognition typographical errors in historical OCR texts (line-break splits, hyphens, and isolated broken characters verified against context or the French original) are normalized in display quotations and recorded in the accompanying source audit.

Restricted or catalog-only magisterial sources---specifically the \emph{Catechism of the Catholic Church}, Pope Benedict XII's \emph{Benedictus Deus}, the Council of Trent's Sessions VI and XIII decrees, St.\ John Paul II's \emph{Ecclesia de Eucharistia}, and Pope Benedict XVI's \emph{Spe salvi}---are cited strictly by accurate theological paraphrase with zero verbatim reproduction.

A search of the repository source library confirmed that no works of St.\ John of the Cross and St.\ John Paul II's apostolic letter \emph{Salvifici doloris} (1984) are currently registered there; Pope Benedict XVI's encyclical \emph{Deus caritas est} (2005) is registered as a Vatican catalog entry under papal documents (restricted and unquoted) and lies outside the sources cited here. The spiritual theology of dryness and darkness is carried by directly verified witnesses (St.\ Teresa of Ávila, St.\ Francis de Sales, and St.\ Th\'er\`ese of Lisieux).

Public-domain texts and translations retain their public-domain status; original exposition, structure, and synthesis composed for this essay are licensed under CC BY 4.0.

\subsection*{Review State}

All primary quotations, locus citations, source bindings, schema checks, web-edition conversions, and print builds were verified internally against tracked repository artifacts on 2026-11-28, including an independent cold review and full remediation recorded in the accompanying cold-review log. No external specialist, clerical, or peer review is claimed.
